% SPDX-License-Identifier: Apache-2.0
% covers: FrameLen
\datasheet{FrameLen}{An Ethernet frame held whole and handed on with its
length, so that a receiver's length can be found again after a clock
crossing.}

\dsfacts{\code{lib/parts/src/ethdma.rs}}{\code{txhdl\_parts::ethdma::FrameLen}}%
{\code{//docs:examples}}

\subsection*{Function}

\code{FrameIn} must know a frame's length before the store engine
starts. \code{EthRx} knows it and says so, but on the board the
receiver is in the top, on the clock the PHY recovers, and only the
bytes cross to the core's clock.

The length cannot follow the bytes on a wire of its own, because it
would arrive on the wrong clock and out of step with the frame it
describes. It cannot ride inside the stream either, because the remote
peripheral reads the same stream and would read the length as a byte
of its frame.

So the length is found again after the crossing. \code{FrameLen}
takes a frame a byte at a time, holds all of it, and then offers it
back with its count on \code{len}, which is exactly what \code{EthRx}
offers. It is the receiver's output side without the receiver:
\code{FrameIn} joined to it cannot tell the difference.

\subsection*{Parameters}

None. It holds \code{FRAME\_MAX} bytes, 2048, the same as the
receiver.

\dstables{FrameLen}

\subsection*{Behaviour}

\begin{itemize}
\item It takes bytes until one is marked last, then hands the frame
on, the held bytes in order and the last one marked. It takes no byte
while it is handing one on.
\item \code{len} is the count held while a frame is being handed on,
and zero otherwise, as the receiver's own length is, so that a stale
count cannot be read as a current one.
\end{itemize}

\subsection*{Verification}

\begin{itemize}
\item \code{ex\_framelen} puts it between \code{FrameOut} and
\code{FrameIn}, in the place \code{EthRx} has in \code{ex\_ethdma},
and sends the same two frames, of 101 and 67 bytes. It checks every
word stored, that the length \code{FrameIn} was told for the last frame
is 67, and that the slots alternated.
\item The netlist is checked against that run under nvc and
Verilator: \code{//docs:framelen\_sim\_framelen\_tb\_test} and
\code{//docs:framelen\_vsim\_test}.
\item On the board, it is in the path of
\code{a\_frame\_received\_lands\_in\_memory\_and\_the\_core\_reads\_it\_back}
in \code{cpu/vreteno/tests/board.rs}.
\end{itemize}

\subsection*{Used by}

\code{Board}, as \code{flen}, between \code{EthShare}'s second user
and \code{FrameIn}. The example \code{ex\_framelen}.

\subsection*{Limits}

It costs a frame of memory and a frame of latency, since the frame was
already held whole once, before the crossing. That is the price of the
length not crossing with it.

A frame longer than 2047 bytes is handed on as its first 2047, the
last of them marked. The bytes after those are written over the one
slot of memory that is not handed on, rather than wrapping round onto
the start of the frame.

It runs on one clock, \code{DefaultClock}.
